\section{Conclusion and Future Work}\label{sec:conclusion}

This study presented a controlled experimental framework for evaluating
language reasoning, retrieval grounding, and hybrid fusion within stock
market forecasting systems.

Four forecasting models were developed and evaluated under identical
experimental conditions. Although the hybrid model achieved the highest
prediction accuracy, statistical analysis indicated that the observed
improvements were not significant.

The experiments revealed that retrieval techniques improve response
quality and contextual grounding, whereas hybrid modelling enhances
consistency and interpretability. Nevertheless, prediction variability
and numerical errors remain important challenges.

The principal contribution of this work is the introduction of a
reproducible evaluation framework capable of measuring the individual
impact of language reasoning, retrieval mechanisms, and model fusion
techniques.

Future work may include larger datasets, more sophisticated retrieval
methods, advanced fusion strategies, and alternative language-model
architectures.

\subsection{Future Work}\label{sec:future-work}

\begin{itemize}
\item
  \textbf{Power where the effect is.} The most valuable extension is not
  more stocks but more shock. The study\textquotesingle s strongest
  effects all sit in the 110 shock rows, and the significance question
  turns entirely on that count. Extending the period back through
  2008-2012 would roughly triple the shock sample without changing any
  other part of the design, and would settle whether the shock-regime
  gap is real.
\item
  \textbf{Isolate the unit-error mechanism.} The unit-error finding is
  the most transferable result here, and it was found incidentally. A
  study that varied the amount of retrieved context - two items, four,
  eight - and measured numerical fidelity at each level would establish
  whether the effect scales with context length, which would matter to
  any retrieval-augmented system that quotes figures.
\item
  \textbf{Let fusion write its own answer.} The fusion inherits the
  grounded arm\textquotesingle s text unchanged, which is why it cannot
  beat it on grounding or helpfulness. A version that regenerates the
  explanation conditioned on both arms\textquotesingle{} calibrated
  scores would be able to, and would test whether the 0.8-point fusion
  gain is a ceiling or an artefact of the current design.
\item
  \textbf{Validate the rubric against people.} The rubric scores what
  the system delivers; it does not establish that a user acts better as
  a result. A small study in which participants make decisions with each
  arm\textquotesingle s output would test whether the 26.6-point spread
  corresponds to anything a person experiences.
\end{itemize}

\nocite{*}
